\documentclass[11pt,a4paper]{article}
\usepackage[utf8]{inputenc}
\usepackage[T1]{fontenc}
\usepackage{geometry}
\geometry{margin=1in}
\usepackage{hyperref}
\usepackage{enumitem}
\usepackage{booktabs}
\usepackage{longtable}
\usepackage{xcolor}
\usepackage{titlesec}
\usepackage{amsmath,amssymb}

\titleformat{\section}{\Large\bfseries\color{blue!70!black}}{}{0em}{}
\titleformat{\subsection}{\large\bfseries\color{blue!50!black}}{}{0em}{}
\titleformat{\subsubsection}{\normalsize\bfseries\color{blue!30!black}}{}{0em}{}

\title{\textbf{Replication Plan}\\[0.5em]
\large FDTD: Solving 1+1D Delay PDE in Parallel\\[0.3em]
\normalsize OSTI ID: 1475143}
\author{Auto-generated Replication Blueprint}
\date{\today}

\begin{document}
\maketitle
\tableofcontents
\newpage

%---------------------------------------------------------------------
\section{Paper Overview}
%---------------------------------------------------------------------
This paper presents a finite-difference time-domain (FDTD) method for solving 1+1D partial differential equations with time delays. The authors develop a parallelization strategy that exploits the causal structure of delay PDEs---the delay term means that future values depend on past values at a fixed lag, enabling a pipeline-parallel approach. The paper demonstrates the method on specific delay PDEs (e.g., delay diffusion equation, delay wave equation), analyzes stability (CFL-like conditions modified for the delay), and benchmarks parallel speedup.

%---------------------------------------------------------------------
\section{Detailed Workflow}
%---------------------------------------------------------------------

\subsection{Phase 1: FDTD Scheme Implementation}
\begin{enumerate}[label=\textbf{1.\arabic*}]
  \item \textbf{Discretize the delay PDE.}
    \begin{itemize}
      \item Spatial domain: $x \in [0, L]$ with $N_x$ grid points, spacing $\Delta x = L / N_x$.
      \item Temporal domain: $t \in [0, T]$ with $N_t$ time steps, spacing $\Delta t$.
      \item Delay: $\tau > 0$ (fixed), requiring storage of solution at times $t - \tau$.
    \end{itemize}
  \item \textbf{Implement the FDTD update rule.}
    \begin{itemize}
      \item Standard central differences for spatial derivatives.
      \item Forward Euler or Crank--Nicolson for time stepping.
      \item Delay term: look up $u(x, t - \tau)$ from stored history buffer.
      \item Implement boundary conditions (Dirichlet, Neumann, periodic as specified).
    \end{itemize}
  \item \textbf{History buffer management.}
    \begin{itemize}
      \item Store solution snapshots for $t \in [t_n - \tau, t_n]$; use circular buffer.
      \item Handle interpolation if $\tau / \Delta t$ is not an integer.
    \end{itemize}
\end{enumerate}

\subsection{Phase 2: Stability Analysis}
\begin{enumerate}[label=\textbf{2.\arabic*}]
  \item \textbf{Derive and verify CFL-like stability condition} for the delay FDTD scheme.
  \item \textbf{Numerically test stability boundaries.}
    \begin{itemize}
      \item Run simulations at various $\Delta t / \Delta x$ ratios near the stability boundary.
      \item Verify that solutions blow up beyond the predicted threshold.
    \end{itemize}
\end{enumerate}

\subsection{Phase 3: Serial Verification}
\begin{enumerate}[label=\textbf{3.\arabic*}]
  \item \textbf{Solve test problems with known analytical solutions.}
    \begin{itemize}
      \item Delay diffusion equation with manufactured solution.
      \item Verify convergence order (second-order in space, first/second in time).
    \end{itemize}
  \item \textbf{Reproduce solution plots} from the paper (space-time contour plots, time traces at fixed $x$).
\end{enumerate}

\subsection{Phase 4: Parallelization}
\begin{enumerate}[label=\textbf{4.\arabic*}]
  \item \textbf{Implement pipeline-parallel FDTD.}
    \begin{itemize}
      \item Partition the time domain into blocks of size $\tau$.
      \item Within each block, the delay term references only the previous block (already computed).
      \item Use MPI or multiprocessing to assign blocks to different processors.
    \end{itemize}
  \item \textbf{Implement spatial-domain decomposition} (standard for FDTD) as a comparison.
  \item \textbf{Implement hybrid spatial--temporal parallelism} if described.
\end{enumerate}

\subsection{Phase 5: Performance Benchmarking}
\begin{enumerate}[label=\textbf{5.\arabic*}]
  \item \textbf{Measure wall-clock time} vs.\ number of processors for fixed problem size (strong scaling).
  \item \textbf{Measure efficiency} as problem size grows with fixed processors (weak scaling).
  \item \textbf{Reproduce speedup plots} from the paper.
  \item \textbf{Compare parallel vs.\ serial solutions} for numerical equivalence (bitwise or within tolerance).
\end{enumerate}

%---------------------------------------------------------------------
\section{Datasets}
%---------------------------------------------------------------------
All data is synthetically generated by the FDTD solver. No external datasets are needed.

\begin{longtable}{p{4.5cm} p{3.5cm} p{6cm}}
\toprule
\textbf{Dataset / Source} & \textbf{Type} & \textbf{Location / Notes} \\
\midrule
\endfirsthead
\toprule
\textbf{Dataset / Source} & \textbf{Type} & \textbf{Location / Notes} \\
\midrule
\endhead
Initial/boundary conditions & Synthetic & Specified analytically in the paper \\
\midrule
Manufactured solutions & Analytic & Derived from the PDE for convergence tests \\
\midrule
Timing benchmarks & Generated & Collected from parallel runs on the target hardware \\
\bottomrule
\end{longtable}

%---------------------------------------------------------------------
\section{Models, Tools, and Software}
%---------------------------------------------------------------------

\subsection{Programming Languages}
\begin{itemize}
  \item \textbf{Python 3.8+} --- for prototyping and serial implementation.
  \item \textbf{C/C++ or Fortran} --- for performance-critical parallel implementation.
\end{itemize}

\subsection{Libraries and Frameworks}
\begin{longtable}{p{4.5cm} p{2.5cm} p{6.5cm}}
\toprule
\textbf{Tool / Library} & \textbf{Version} & \textbf{Purpose / URL} \\
\midrule
\endfirsthead
\toprule
\textbf{Tool / Library} & \textbf{Version} & \textbf{Purpose / URL} \\
\midrule
\endhead
NumPy & $\geq 1.21$ & Array operations for serial FDTD \\
\midrule
SciPy & $\geq 1.7$ & Sparse matrices, reference ODE solvers for validation \\
\midrule
mpi4py & $\geq 3.1$ & MPI bindings for Python parallel implementation \newline \url{https://mpi4py.readthedocs.io/} \\
\midrule
MPI (OpenMPI / MPICH) & system & Message passing for distributed parallelism \\
\midrule
Matplotlib & $\geq 3.4$ & Solution visualization, speedup plots \\
\midrule
Numba or Cython & latest & Optional JIT compilation for Python FDTD loops \\
\midrule
multiprocessing & stdlib & Shared-memory parallelism alternative \\
\bottomrule
\end{longtable}

\subsection{Hardware Requirements}
\begin{itemize}
  \item Multi-core workstation or small cluster for parallel benchmarking (4--64 cores).
  \item RAM: modest ($\sim$4\,GB per process) for 1+1D problems.
\end{itemize}

\end{document}
